% !TEX root = ../main.tex
\chapter{\texttt{simp} と \texttt{rw}}
\label{chap:ch08_simp_rw}
\BookIndex{simp}{\texttt{simp}}
\BookIndex{rw}{\texttt{rw}}
\BookIndex{unfold}{\texttt{unfold}}

\begin{chapterabstract}
本章では、\texttt{rfl} の次の証明道具として \texttt{simp}、\texttt{unfold}、\texttt{rw} を導入します。
\texttt{simp} は、Lean が知っている単純化規則で式を整理します。
\texttt{unfold} は、指定した定義を現在のゴールで明示的に展開します。
\texttt{rw} は、すでに証明済みの等式を使って式を書き換えます。
いずれも \texttt{rfl} だけでは届かない場面で使う、等式推論の基本道具です。
\end{chapterabstract}

\begin{learninggoals}
\item \texttt{simp} を使って、分岐や明らかな整理を含む等式を証明できます。
\item \texttt{unfold} を、指定した定義をゴール内で展開するタクティクとして読めます。
\item \texttt{rw} を使って、既知の等式でゴールを書き換えられます。
\item \texttt{rfl}、\texttt{simp}、\texttt{unfold}、\texttt{rw} の役割を読み分けられます。
\end{learninggoals}

\section{定義展開だけでは閉じない証明}

これまでの証明は、\texttt{rfl} だけで閉じてきました。
しかし、分岐を整理する場面や、別の場所で証明した等式を使う場面では、\texttt{rfl} だけでは足りません。
本章では、分岐の整理には \texttt{simp}、定義を明示的に開くには \texttt{unfold}、証明済みの等式による書き換えには \texttt{rw} を導入します。
これらは、等式推論とリファクタリング証明で中心的に使います。

\section{\texttt{simp} で明らかな整理を任せる}

\texttt{simp} は、Lean が知っている単純化規則を使って式を整理する道具です。
たとえば、条件が \texttt{true} に決まっている分岐を整理できます。

\begin{listing}[htbp]
\begin{minted}{lean4}
def chargeShipping (isMember : Bool) (shipping : Nat) : Nat :=
  if isMember then 0 else shipping

theorem member_shipping_zero (shipping : Nat) :
    chargeShipping true shipping = 0 := by
  simp [chargeShipping]

-- これは任意の shipping については証明できない。
-- shipping が 0 とは限らないため。
-- theorem nonmember_shipping_zero (shipping : Nat) :
--     chargeShipping false shipping = 0 := by
--   simp [chargeShipping]
\end{minted}
\caption{\texttt{member\_shipping\_zero}}
\label{lst:ch08_simp}
\end{listing}

コード\ref{lst:ch08_simp}では、\texttt{chargeShipping true shipping} は \texttt{if true then 0 else shipping} になり、\texttt{simp [chargeShipping]} がこの式を \texttt{0} に整理します。
角括弧の \texttt{[chargeShipping]} は、「この定義も使って整理してよい」という指定です。
コメントアウトした \texttt{nonmember\_shipping\_zero} は、非会員の場合に \texttt{shipping} がそのまま残るため、任意の \texttt{shipping} については証明できません。

\begin{table}[htbp]
\centering
\scriptsize
\setlength{\tabcolsep}{3pt}
\caption{\texttt{chargeShipping} の二つの分岐}
\label{tab:ch08-simp-branches}
\begin{tabular}{@{}p{0.20\linewidth}p{0.74\linewidth}@{}}
\toprule
式 & 展開後 \\
\midrule
\texttt{chargeShipping true shipping} & \texttt{0} \\
\texttt{chargeShipping false shipping} & \texttt{shipping} \\
\bottomrule
\end{tabular}
\end{table}

\section{\texttt{unfold} で定義を明示的に展開する}

\texttt{unfold} は、\texttt{def} や \texttt{theorem} のように新しい宣言を作る語ではなく、\texttt{by} 以下で現在の証明ゴールを変形する\emph{タクティク}です。
\texttt{unfold chargeShipping} と書くと、ゴール内に現れる \texttt{chargeShipping} を、その定義本体へ一段展開します。
定義の中身を先に見える形へしてから、別のタクティクで残りを処理したいときに使います。

前節の \texttt{simp [chargeShipping]} は、\texttt{chargeShipping} の展開と、その後の単純化をまとめて行います。
同じ流れを分けて書く場合は、まず \texttt{unfold chargeShipping} で定義を展開し、続けて \texttt{simp} で \texttt{if true then 0 else shipping} を整理します。
前者は短く書きたい場合、後者は「定義展開」と「単純化」を証明手順として明示したい場合に適しています。

\begin{table}[htbp]
\centering
\caption{\texttt{unfold} と \texttt{simp} の役割}
\label{tab:ch08-unfold-simp}
\begin{tabular}{p{0.34\linewidth}p{0.56\linewidth}}
\toprule
操作 & 役割 \\
\midrule
\texttt{unfold chargeShipping} & 指定した定義の本体を現在のゴールに現します。 \\
\texttt{simp} & 展開後の分岐や既知の単純化規則を整理します。 \\
\texttt{simp [chargeShipping]} & 指定した定義の展開と単純化をまとめて行います。 \\
\bottomrule
\end{tabular}
\end{table}

複数の定義を続けて展開する場合は、名前を空白で並べます。
次章に現れる \texttt{unfold rawCalc refactoredCalc} は、現在のゴールにある \texttt{rawCalc} と \texttt{refactoredCalc} の二つを展開する指定です。

\section{\texttt{rw} で既知の等式を使う}

\texttt{rw} は rewrite の略であり、すでに証明済みの等式を使って式の一部を書き換えます。
書き換えた結果、左右が同じ形になれば証明が閉じます。

\begin{listing}[htbp]
\begin{minted}{lean4}
def addTax (price tax : Nat) : Nat :=
  price + tax

def normalizePrice (price : Nat) : Nat :=
  price

theorem normalizePrice_eq (price : Nat) :
    normalizePrice price = price := by
  rfl

theorem addTax_after_normalize (price tax : Nat) :
    addTax (normalizePrice price) tax = addTax price tax := by
  rw [normalizePrice_eq]
\end{minted}
\caption{\texttt{addTax\_after\_normalize}}
\label{lst:ch08_rw}
\end{listing}

コード\ref{lst:ch08_rw}では、\texttt{normalizePrice\_eq} という等式を使って、\texttt{normalizePrice price} を \texttt{price} に書き換えています。
書き換えの過程は、次のように追えます。

\begin{table}[htbp]
\centering
\caption{\texttt{normalizePrice\_eq} による書き換え}
\label{tab:ch08-rw}
\begin{tabular}{ll}
\toprule
段階 & ゴール \\
\midrule
最初のゴール & \path{addTax (normalizePrice price) tax = addTax price tax} \\
\path{normalizePrice price} を書き換え & \path{addTax price tax = addTax price tax} \\
結論 & 左右が同じなので閉じる \\
\bottomrule
\end{tabular}
\end{table}

\texttt{rfl} は「定義を見れば同じ」、\texttt{simp} は「明らかな整理をまとめて行う」、\texttt{unfold} は「指定した定義の本体をゴールに現す」、\texttt{rw} は「既知の等式で置き換える」と読みます。
これらが、等式推論とリファクタリング証明の基本語彙になります。

\section{本章のまとめ}

要点と記法
\begin{itemize}
\item \texttt{simp}：既知の単純化規則による整理
\item \texttt{simp [name]}：指定した定義や補題も使う整理
\item \texttt{unfold name}：指定した定義の本体を現在のゴールに現す
\item \texttt{unfold f g}：複数の定義を現在のゴールで順に展開する
\item \texttt{rw [h]}：等式 \texttt{h} によるゴールの書き換え
\end{itemize}
